\section{Data Acquisition, Conversion and Distribution}\label{sec:acq}
\lecdate{23 September 2019}

\subsection{Quantizing and quantization error}\scanref{pp.\,1}
The main function of an A/D conversion consists of three steps.
\begin{itemize}
  \item \textbf{Sampling:} the analog signal becomes a discrete-time signal (real numbers).
  \item \textbf{Amplitude quantization:} each number is replaced by an approximation taken from a finite
        set of discrete values. This causes the \emph{quantization error}.
  \item \textbf{Coding:} the discrete values are converted to binary values.
\end{itemize}

\paragraph{Quantization level.} With $\mathrm{FSR}$ the full-scale range and $n$ the number of bits,
\begin{equation}
  Q=\frac{\mathrm{FSR}}{2^{n}} .
\end{equation}

\begin{example}
For an input range of $\pm5\,\mathrm{V}$ (so $\mathrm{FSR}=10\,\mathrm{V}$) and a $12$-bit A/D converter (ADC),
\[
  Q=\frac{10}{2^{12}}=2.44\ \mathrm{mV}.
\]
This ADC cannot generate a new discrete value if the voltage change between two adjacent sample values is
smaller than $2.44\ \mathrm{mV}$.
\end{example}

\paragraph{Quantization error.} Let $x(t)$ be the analog input and $y(t)$ the closest available digital
value at the output of the A/D block. The error is $e(t)=x(t)-y(t)$. It varies between $0$ and $\mp\tfrac12 Q$,
so its magnitude satisfies
\begin{equation}
  0\le \lvert e(t)\rvert\le \tfrac12 Q .
\end{equation}

\begin{center}
\begin{tikzpicture}[scale=1]
  \begin{axis}[width=0.78\linewidth,height=5.6cm,axis lines=left,xlabel={$x(t)$ (analog input)},
      ylabel={$y(t)$ (digital value)},xmin=0,xmax=8,ymin=0,ymax=8.6,
      xtick={0,1,...,8},ytick={0,1,...,8},
      xticklabels={0,,,,,,,,},yticklabels={0,,,,,,,,},grid=both,
      grid style={very thin,gridgray},every axis plot/.append style={thick}]
    \addplot[itublue,const plot mark mid,domain=0:8,samples=17,unbounded coords=jump]
      {round(x)};
    \addplot[itugold,dashed,domain=0:8] {x};
    \node[font=\small,itublue,anchor=west] at (axis cs:4.7,2.5) {$y=$ quantized};
    \node[font=\small,itugold,anchor=west] at (axis cs:1.1,5.4) {$y=x$ (ideal)};
    \draw[<->,boardgreen,thick] (axis cs:5.2,5.2) -- node[right,font=\scriptsize]{$Q$} (axis cs:5.2,6.2);
  \end{axis}
\end{tikzpicture}\\[2pt]
{\small Staircase characteristic of a quantizer; the step is $Q$ and the error never exceeds $Q/2$.}
\end{center}

\begin{addsol}
\textbf{Problem.} Find $Q$ for a $\pm5\,\mathrm{V}$ converter with $n=8,10,12,16$ bits, and the number of
levels $2^n$.

\medskip
\begin{center}
\begin{tabular}{@{}rrr@{}}\toprule
$n$ & $2^n$ & $Q=10/2^n$ \\ \midrule
8  & 256    & 39.06 mV \\
10 & 1024   & 9.77 mV \\
12 & 4096   & 2.44 mV \\
16 & 65536  & 0.153 mV \\ \bottomrule
\end{tabular}
\end{center}
The maximum error is $Q/2$, e.g.\ $1.22\,\mathrm{mV}$ for $n=12$.
\end{addsol}

%------------------------------------------------------------------------------
\subsection{Data acquisition, conversion and distribution systems}\scanref{pp.\,2--6}
The signal conversion in a digital control system involves the following operations:
\begin{enumerate}
  \item Multiplexing and demultiplexing,
  \item Sample and hold,
  \item A/D conversion (quantizing, encoding),
  \item D/A conversion (decoding).
\end{enumerate}

\paragraph{Data-acquisition system.} Physical variables (position, acceleration, temperature, pressure,
\ldots) are converted to a voltage or current by a transducer, amplified, low-pass filtered (noise
cancellation), multiplexed, sampled-and-held and finally converted by the A/D converter before they go to
the digital controller.

\begin{center}
\resizebox{\linewidth}{!}{\begin{tikzpicture}[node distance=13mm,font=\small]
  \node[blk,minimum height=9mm] (tr) {Transducer};
  \node[blk,right=of tr] (am) {Amplifier};
  \node[blk,right=of am,text width=17mm] (lp) {Low Pass Filter};
  \node[blk,right=of lp] (mux) {Analog MUX};
  \node[blk,right=of mux] (sh) {Sample Hold};
  \draw[sig] ($(tr.west)+(-8mm,0)$) -- node[above,pos=0.2,text width=13mm,align=center,font=\scriptsize]{physical variables} (tr);
  \draw[sig] (tr) -- node[above,font=\scriptsize,align=center]{voltage\\ or current} (am);
  \draw[sig] (am) -- (lp);
  \draw[sig] (lp) -- node[below,font=\scriptsize]{noise cancel.} (mux);
  \draw[sig] (mux) -- (sh);
  \node[blk,below=9mm of sh] (adc) {A/D Conv.};
  \draw[sig] (sh) -- (adc);
  \draw[sig] (adc.west) -- ++(-8mm,0) node[left,font=\scriptsize,align=center]{to digital\\control};
  \draw[sig] ($(mux.north)+(0,5mm)$) node[above,font=\scriptsize]{Analog} -- (mux.north);
  \node[font=\scriptsize,right=1mm of adc] {Digital};
\end{tikzpicture}}
\end{center}

\paragraph{Data-distribution system.} The digital value from the controller is stored in a register,
demultiplexed, converted by D/A converters and held before it reaches the actuator.

\begin{center}
\resizebox{\linewidth}{!}{\begin{tikzpicture}[node distance=13mm,font=\small]
  \node[blk] (rg) {Register};
  \node[blk,right=of rg] (dm) {DeMUX};
  \node[blk,right=of dm] (da) {D/A conv.};
  \node[blk,right=of da] (hd) {Hold};
  \draw[sig] ($(rg.west)+(-10mm,0)$) -- node[above,pos=0.1,font=\scriptsize,align=center]{from digital\\control} (rg);
  \draw[sig] (rg) -- (dm);
  \draw[sig] (dm) -- node[above,font=\scriptsize]{Digital} (da);
  \draw[sig] (da) -- (hd);
  \draw[sig] (hd) -- node[above,font=\scriptsize,pos=0.5]{Analog} ++(1.8,0) node[right,font=\scriptsize,align=left]{to\\actuator};
\end{tikzpicture}}
\end{center}

%------------------------------------------------------------------------------
\subsubsection*{Analog multiplexer (MUX)}
\begin{itemize}
  \item It is a multiple switch (usually an electronic switch) that sequentially switches among many analog
        input channels.
  \item Number of channels: $4,8,16,32,64,\dots$
  \item At a given instant only one switch is in the ``on'' position.
\end{itemize}

\begin{center}
\begin{tikzpicture}[x=1cm,y=0.55cm]
  \draw[thick] (0,-0.7) rectangle (3.6,8.6);
  \foreach \i in {0,...,7} {
    \pgfmathsetmacro{\yy}{7.8-\i}
    \draw (-1.2,\yy) -- ++(0.3,0) -- ++(0.1,0.18) -- ++(0.1,-0.36) -- ++(0.1,0.36) -- ++(0.1,-0.36) -- ++(0.1,0.18) -- (0.2,\yy);
    \node[circle,fill,inner sep=1pt] at (0.2,\yy) {};
    \node[circle,draw,inner sep=1pt] at (2.2,\yy) {};
    \draw (3.6,\yy) -- (4.2,\yy);
    \ifnum\i=0
      \draw[thick] (0.2,\yy) -- (2.2,\yy);
    \else
      \draw[thick] (0.2,\yy) -- ($(0.2,\yy)+(1.7,0.45)$);
    \fi
  }
  \draw[thick] (4.2,0.8) -- (4.2,7.8);
  \draw[->] (4.2,4.3) -- (4.9,4.3) node[right,font=\small]{to S/H};
  \node[font=\small,anchor=east] at (-1.3,4.3) {Input channels};
  \draw[->,thick] (1.7,-1.6) -- (1.7,-0.7);
  \node[font=\small] at (1.7,-2.1) {sequencer};
\end{tikzpicture}
\end{center}

\subsubsection*{Demultiplexer}
It separates the composite output digital data from the digital controller into the original channels. Each
channel is connected to a D/A converter to produce the analog signal of that channel.

\begin{center}
\begin{tikzpicture}[font=\small]
  \node[blk,minimum width=10mm] (d1) at (0,0.7) {DAC};
  \node[blk,minimum width=10mm] (d2) at (0,-0.1) {DAC};
  \draw[sig] (-1.6,0.7) -- (d1); \draw[sig] (-1.6,-0.1) -- (d2);
  \draw[thick] (-1.6,0.7) -- (-1.6,-1.1);
  \draw[sig] (-2.2,-0.2) -- (-1.6,-0.2);
  \draw[sig] (-1.6,-0.9) -- (-1.1,-0.9) node[right]{\vdots};
  \draw[sig] (d1) -- ++(1.2,0) node[right]{Analog 1};
  \draw[sig] (d2) -- ++(1.2,0) node[right]{Analog 2};
\end{tikzpicture}
\end{center}

\subsubsection*{Sample-and-hold circuit}
A sampler converts an analog signal into a train of amplitude-modulated pulses. The hold circuit holds the
value of the sampled pulse over a specified period of time. The S/H circuit is an analog circuit (simply a
voltage-memory device).

\begin{center}
\begin{tikzpicture}
  \begin{axis}[width=0.8\linewidth,height=5.2cm,axis lines=left,xlabel={$t$},ylabel={$x$},
      xmin=0,xmax=10.5,ymin=0,ymax=1.35,xtick=\empty,ytick=\empty,clip=false]
    \addplot[itublue,thick,domain=0:10,samples=100,smooth]
      {0.3+0.95*sin(deg(x*0.33))^1.2*exp(-0.04*x)};
    \addplot[ycomb,black!70,thick,samples at={0.5,1.5,...,9.5}]
      {0.3+0.95*sin(deg(x*0.33))^1.2*exp(-0.04*x)};
    \addplot[const plot,itugold,thick,samples at={0.5,1.5,...,9.5}]
      {0.3+0.95*sin(deg(x*0.33))^1.2*exp(-0.04*x)};
  \end{axis}
\end{tikzpicture}\\[2pt]
{\small Blue: analog signal; vertical lines: sampled pulses; gold: held (staircase) value.}
\end{center}

%------------------------------------------------------------------------------
\subsection{The AD585 sample-and-hold IC and types of A/D converters}
IC AD585 (Analog Devices) is a high-speed, precision S/H circuit: a unity-gain buffer (Op-Amp~1), a switch
operated by the sample/hold command, a hold capacitor and an output buffer (Op-Amp~2).

\begin{center}
\resizebox{\linewidth}{!}{% AD585 sample-and-hold (cf. handwritten page 4)
\begin{tikzpicture}[font=\small,line cap=round]
  % buffer 1
  \node[isosceles triangle,draw=itublue,thick,minimum height=13mm,inner sep=0pt,fill=white] (b1) at (1.6,0) {};
  \draw (b1.east) coordinate(o1);
  \draw[thick] (-0.3,0) node[left]{Analog input} to[short,o-] (b1.west);
  \draw[thick] (o1) -- ++(0.35,0) coordinate(j1) |- ($(b1.north west)+(0,0.45)$) -- ($(b1.north west)+(0,0.45)$) -- ($(b1.north west)!0.5!(b1.south west)+(0,0.3)$);
  \draw[thick] (o1) -- (3.5,0);
  % switch
  \draw[thick] (3.5,0) -- (3.9,0);
  \draw[thick] (3.9,0) -- (4.5,0.28);
  \draw[thick] (4.7,0) -- (5.3,0) coordinate(hold);
  \node[circle,draw,inner sep=1pt] at (3.9,0) {}; \node[circle,draw,inner sep=1pt] at (4.7,0) {};
  % hold capacitor to ground
  \draw[thick] (5.3,0) -- (5.3,-0.75);
  \draw[thick] (5.05,-0.75) -- (5.55,-0.75); \draw[thick] (5.05,-1.0) -- (5.55,-1.0);
  \draw[thick] (5.3,-1.0) -- (5.3,-1.6) node[ground]{};
  \node[font=\scriptsize,align=left,anchor=west] at (5.55,-0.9) {hold\\capacitor};
  % buffer 2
  \node[isosceles triangle,draw=itublue,thick,minimum height=13mm,inner sep=0pt,fill=white] (b2) at (7.2,0) {};
  \draw[thick] (5.3,0) -- (b2.west);
  \draw[thick] (b2.east) -- ++(0.35,0) coordinate(j2) |- ($(b2.north west)+(0,0.45)$) -- ($(b2.north west)+(0,0.45)$) -- ($(b2.north west)!0.5!(b2.south west)+(0,0.3)$);
  \draw[thick] (b2.east) -- (8.9,0) node[circle,draw,inner sep=1pt]{} node[right=2mm,align=left]{Analog\\output};
  % ground line
  \draw[thick] (-0.3,-2.2) node[circle,draw,inner sep=1pt]{} -- (8.9,-2.2) node[circle,draw,inner sep=1pt]{};
  \draw[thick] (5.3,-1.6) -- (5.3,-2.2);
  % command
  \node[draw,fill=white,font=\scriptsize] (cmd) at (4.3,-1.5) {Sampler};
  \node[font=\scriptsize,anchor=east,align=right] at (cmd.west) {Sample and hold\\command};
  \draw[dashed,->] (cmd.north) -- (4.3,0.05);
  \node[font=\scriptsize] at (1.6,-1.0) {Op-Amp 1}; \node[font=\scriptsize] at (7.2,-1.0) {Op-Amp 2};
\end{tikzpicture}
}
\end{center}

\paragraph{Types of A/D converters.} The following types are used most frequently: successive-\allowbreak approximation, integrating, counter and parallel (flash) type.\par

\begin{center}
\begin{tabular}{@{}clccc@{}}\toprule
Type & Name & Speed & Resolution & Noise immunity\ \textbar\ Cost \\ \midrule
1 & Successive approximation & medium & 10--16 bits & poor \ \textbar\ low \\
2 & Integrating & slow & 12--18 bits & good \ \textbar\ low \\
3 & Counter & slow & 14--24 bits & good \ \textbar\ medium \\
4 & Parallel / flash & fast & 4--8 bits & none \ \textbar\ high \\ \bottomrule
\end{tabular}
\end{center}

\paragraph{Errors in A/D converters.} Actual converters show \emph{offset error}, \emph{linearity error} and
\emph{gain error} with respect to the ideal characteristic (binary code against input, from $000$ up to
$111$ at full scale FS).

\begin{center}
\resizebox{\linewidth}{!}{% A/D converter errors (cf. handwritten page 4)
\begin{tikzpicture}[font=\scriptsize]
  \def\W{3.0}\def\H{2.6}
  \foreach \k in {0,4.0,8.0} {
    \begin{scope}[xshift=\k cm]
      \draw[->] (0,0) -- (\W+0.3,0); \draw[->] (0,0) -- (0,\H+0.4) node[above]{BIN};
      \node[left] at (0,0) {000}; \node[left] at (0,\H/2) {100}; \node[left] at (0,\H) {111};
      \draw (\W/2,-0.05) -- (\W/2,0.05); \draw (\W,-0.05) -- (\W,0.05);
      \node[below] at (\W/2,0) {$\tfrac12$FS}; \node[below] at (\W,0) {FS};
      \draw[dashed,black!70] (0,0) -- (\W,\H);
    \end{scope}
  }
  % panel 1: offset error
  \draw[itublue,thick] (0.55,0) -- (\W,{\H*(\W-0.55)/\W});
  \node[black!70,rotate=40] at (1.0,1.65) {ideal};
  \node[itublue] at (2.15,1.45) {actual};
  \draw[<-] (1.35,0.5) -- (1.9,0.2) node[right]{offset error};
  % panel 2: linearity error
  \begin{scope}[xshift=4.0cm]
    \draw[itublue,thick,domain=0:\W,samples=40,smooth] plot (\x,{\H*(1-exp(-2.2*\x/\W))/(1-exp(-2.2))});
    \node[itublue] at (1.1,2.35) {actual}; \node[black!70] at (2.5,1.55) {ideal};
    \draw[<-] (1.3,1.35) -- (1.9,0.65) node[right]{linearity};
    \node at (2.45,0.42) {error};
  \end{scope}
  % panel 3: gain error
  \begin{scope}[xshift=8.0cm]
    \draw[itublue,thick] (0,0) -- (\W*0.62,\H);
    \node[itublue] at (0.75,2.5) {actual}; \node[black!70] at (2.4,1.35) {ideal};
    \draw[<-] (1.45,1.45) -- (1.95,1.0) node[right]{gain error};
  \end{scope}
\end{tikzpicture}
}
\end{center}

%------------------------------------------------------------------------------
\subsection{Types of D/A converters}\scanref{pp.\,5}
For a full-range digital input there are $2^{n}$ corresponding different analog values, including $0$. For
example, an $8$-bit DAC has $2^{8}=256$ values, i.e.\ the codes $0$--$255$. In D/A conversion there is a
\emph{one-to-one correspondence} (so there is no quantization error).

Two methods are used: \textbf{weighted resistors} (simple but the accuracy is low) and the \textbf{R-2R
ladder network} (complicated but the accuracy is high).

\paragraph{4-bit DAC with weighted resistors.} The inverting op-amp sums the currents of the branches
$8R,4R,2R,R$; $b_1$ is the LSB and $b_4$ the MSB.
\begin{equation}
  V_{\mathrm{out}}=\frac{R_0}{R}\Bigl[b_4+\frac{b_3}{2}+\frac{b_2}{4}+\frac{b_1}{8}\Bigr]V_{\mathrm{REF}} .
\end{equation}

\begin{center}
\resizebox{\linewidth}{!}{% 4-bit weighted-resistor DAC (cf. handwritten page 5)
\begin{tikzpicture}[european resistors,font=\small]
  \foreach \i/\r/\b/\x in {1/8R/b_1/0, 2/4R/b_2/1.5, 3/2R/b_3/3.0, 4/R/b_4/4.5} {
    \draw (\x,2.0) to[R,l=\r,i_=] (\x,0.4);
    \draw (\x,0.4) -- (\x,-0.2);
    \node[circle,draw,inner sep=1pt] at (\x,-0.3) {};
    \draw[thick] (\x,-0.3) -- (\x+0.42,-0.95) node[right,font=\scriptsize]{$\b$};
    \draw (\x,-1.1) -- (\x,-2.3);
    \node[circle,fill,inner sep=1pt] at (\x,-2.3) {};
  }
  \draw (-0.8,-2.3) node[left]{$V_{\mathrm{REF}}$} -- (4.5,-2.3);
  \draw (0,2.0) -- (6.6,2.0);
  % op-amp (inverting) drawn from basic shapes
  \draw[thick,fill=white,draw=itublue] (6.6,3.0) -- (6.6,1.0) -- (8.8,2.0) -- cycle;
  \node[font=\scriptsize] at (6.9,2.55) {$-$}; \node[font=\scriptsize] at (6.9,1.45) {$+$};
  \draw (6.6,1.45) -- (6.1,1.45) -- (6.1,0.9) node[ground]{};
  \draw (6.6,2.0) -- (6.6,2.0);
  \draw (6.0,2.0) -- (6.0,3.7) to[R,l=$R_0$] (9.4,3.7) -- (9.4,2.0) -- (8.8,2.0);
  \draw (6.0,2.0) node[circle,fill,inner sep=1pt]{};
  \draw (9.4,2.0) node[circle,fill,inner sep=1pt]{} -- (10.0,2.0) node[circle,draw,inner sep=1pt]{} node[right]{$V_{\mathrm{out}}$};
  \node[font=\scriptsize] at (0,-0.8) {LSB}; \node[font=\scriptsize] at (4.5,-3.0) {MSB};
  \node[font=\scriptsize,align=left,fill=white,draw,inner sep=3pt] at (8.2,-1.2)
    {4-bit D/A converter\\using weighted resistors};
\end{tikzpicture}
}
\end{center}

\begin{boardnote}
The op-amp is an \emph{inverting} amplifier, so the output voltage actually carries a minus sign. The board
formula gives the magnitude.
\end{boardnote}

\paragraph{$n$-bit DAC with an R-2R ladder.}
\begin{equation}
  V_{\mathrm{out}}=\frac{V_{\mathrm{REF}}}{2}\Bigl[b_{n-1}+\frac{1}{2}b_{n-2}+\cdots+\frac{1}{2^{\,n-1}}b_{0}\Bigr].
\end{equation}

\begin{center}
\resizebox{\linewidth}{!}{% n-bit R-2R ladder DAC (cf. handwritten page 5)
\begin{tikzpicture}[european resistors,font=\small]
  % ladder
  \draw (0,0) to[R,l=$R$] (2.2,0) to[R,l=$R$] (4.4,0);
  \node at (5.2,0) {$\cdots$};
  \draw (6.0,0) to[R,l=$R$] (8.0,0) to[R,l=$R$] (10.0,0);
  \draw (-1.4,0) to[R,l=$2R$] (0,0) -- ++(-1.4,0);
  \draw (-1.4,0) -- (-1.4,-0.6) node[ground]{};
  \foreach \x/\b in {0/b_0,2.2/b_1,4.4/b_2,8.0/b_{n-2},10.0/b_{n-1}} {
    \node[circle,fill,inner sep=1pt] at (\x,0) {};
    \draw (\x,0) to[R,l_=$2R$] (\x,-1.7);
    \node[circle,draw,inner sep=1pt] at (\x,-1.8) {};
    \draw[thick] (\x,-1.8) -- (\x+0.35,-2.55) node[right,font=\scriptsize]{$\b$};
    \draw (\x,-2.7) -- (\x,-3.4); \node[circle,fill,inner sep=1pt] at (\x,-3.4) {};
  }
  \draw (-0.8,-3.4) node[left]{$V_{\mathrm{REF}}$} -- (10.0,-3.4);
  \node[font=\scriptsize] at (0,-4.0) {LSB}; \node[font=\scriptsize] at (10,-4.0) {MSB};
  % op-amp
  \draw (10.0,0) -- (11.4,0);
  \draw[thick,fill=white,draw=itublue] (11.4,0.9) -- (11.4,-0.9) -- (13.4,0) -- cycle;
  \node[font=\scriptsize] at (11.7,0.45) {$-$}; \node[font=\scriptsize] at (11.7,-0.45) {$+$};
  \draw (11.4,-0.45) -- (11.0,-0.45) -- (11.0,-0.9) node[ground]{};
  \draw (11.0,0) -- (11.0,1.4) to[R,l=$3R$] (14.0,1.4) -- (14.0,0) -- (13.4,0);
  \node[circle,fill,inner sep=1pt] at (11.0,0) {}; \node[circle,fill,inner sep=1pt] at (14.0,0) {};
  \draw (14.0,0) -- (14.5,0) node[circle,draw,inner sep=1pt]{} node[right]{$V_{\mathrm{out}}$};
\end{tikzpicture}
}
\end{center}

\begin{addsol}
\textbf{Problem.} A 4-bit weighted-resistor DAC has $R_0=R$ and $V_{\mathrm{REF}}=5\,\mathrm{V}$. Find the
output for the code $b_4b_3b_2b_1=1011$, the output step (resolution) and the full-scale output.

\medskip\noindent
\emph{Solution.} With $R_0/R=1$,
\[
  \lvert V_{\mathrm{out}}\rvert=\Bigl[1+\tfrac{0}{2}+\tfrac14+\tfrac18\Bigr]\cdot 5=1.375\cdot5=6.875\ \mathrm{V}.
\]
The LSB is $b_1$ with weight $\tfrac18$, so one step is $5/8=0.625\ \mathrm{V}$ and the code $1011=11$ gives
$11\cdot0.625=6.875\ \mathrm{V}$, in agreement. Full scale (code $1111$): $\tfrac{15}{8}\cdot5=9.375\ \mathrm{V}$.
\end{addsol}

%------------------------------------------------------------------------------
\subsection{Reconstructing the input signal by hold circuits}\scanref{pp.\,6}
\paragraph{Zero-order hold (ZOH).} The hold circuit produces a staircase waveform. Because of its simplicity
the ZOH is commonly used in digital control systems.

\paragraph{First-order hold (FOH).} It retains the value of the previous sample, as well as the present one,
and predicts, by extrapolation, the next value.

\begin{center}
% ZOH and FOH reconstruction (cf. handwritten page 6)
\pgfplotsset{holdax/.style={width=0.48\linewidth,height=4.4cm,axis lines=left,xmin=0,xmax=8.5,
  ymin=0,ymax=1.3,xtick=\empty,ytick=\empty,clip=false}}
\begin{tikzpicture}
  \begin{groupplot}[group style={group size=2 by 1,horizontal sep=1.1cm},holdax]
    \nextgroupplot[title={Zero-order hold}]
      \addplot[black!55,dotted,very thick,domain=0:8,samples=120] {0.30+0.9*sin(deg(x*0.40))*exp(-0.04*x)};
      \addplot[const plot,itublue,thick,samples at={0,1,...,8}] {0.30+0.9*sin(deg(x*0.40))*exp(-0.04*x)};
      \addplot[ycomb,itublue!50,samples at={0,1,...,8},mark=none] {0.30+0.9*sin(deg(x*0.40))*exp(-0.04*x)};
    \nextgroupplot[title={First-order hold}]
      \addplot[black!55,dotted,very thick,domain=0:8,samples=120] {0.30+0.9*sin(deg(x*0.40))*exp(-0.04*x)};
      % causal FOH: x(kT)+ (x(kT)-x((k-1)T))/T (t-kT), jumps at the samples
      \pgfmathdeclarefunction{xf}{1}{\pgfmathparse{0.30+0.9*sin(deg(#1*0.40))*exp(-0.04*#1)}}
      \foreach \k in {1,...,7} {
        \pgfmathsetmacro{\xa}{0.30+0.9*sin(deg(\k*0.40))*exp(-0.04*\k)}
        \pgfmathsetmacro{\xb}{0.30+0.9*sin(deg((\k-1)*0.40))*exp(-0.04*(\k-1))}
        \addplot[itublue,thick,domain=\k:\k+1,samples=2] {\xa+(\xa-\xb)*(x-\k)};
      }
      \addplot[itublue,thick,domain=0:1,samples=2] {0.30};
  \end{groupplot}
\end{tikzpicture}
\\[2pt]
{\small Dotted: original signal; solid: reconstructed signal.}
\end{center}


